\section{Tree Co-Width vs Tree-width}\label{appendix:tree-co-vs-tree}
In this subsection we prove the following.
\begin{lemma}
    Let $G$ be a graph that admits a tree decomposition with bags of size at most $w\in\mathbb N$. Then $G$ admits a tree co-decomposition of width at most $w+3$. In formulas:
    \[
    \text{tree co-width}(G)\le \text{treewidth}(G)+2
    \]
\end{lemma}
\begin{proof}
    Let $(T,\{B_a\}_{a\in V(T)\cup E(T)})$ be a tree decomposition of $G$ of width at most $w-1$. That is, $T$ is a tree and the collection $\{B_a\}$ is a family of induced subgraphs indexed by the nodes and edges\footnote{We aknowledge that this is not the usual definition of tree-decomposition that appears in the literature however, for our purposes, the equivalence is not hard to see.} of $T$ so that the following criteria hold:\vspace{-.25cm}
    \begin{enumerate}
        \item together the bags $\{B_a\}$ cover $G$; and\vspace{-.25cm}
        \item for each vertex $x\in V(G)$, the indices $a\in V(T)\cup E(T)$ for which $x\in V(B_a)$ form a contiguous subtree of $T$; and\vspace{-.25cm}
        \item for every $a\in V(T)\cup E(T)$, we have $|B_a|\le w$.\vspace{-.25cm}
    \end{enumerate}
    We assume using standard methods that that the maximum degree of $T$ is $3$. We will construct a tree co-decomposition $\Ta:=(T,\{G_a\},\{g_e^t\})$, having the same auxiliary tree as the input tree-decomposition, as follows: For each $a\in V(T)\cup E(T)$, let $V_a$ be the partition of $V(G)$ with the following parts:\vspace{-.25cm}
    \begin{enumerate}
        \item for each $x\in B_a$, we have $\{x\}\in V_a$; and\vspace{-.25cm}
        \item for each connected component $C$ of $T-a$, let $P_C\subseteq V(G)$ be the set of all vertices of $G$ in bags indexed in $C$ that are not in $V(B_a)$, and include $P_C\in V_a$ if $P_C$ is nonempty.\vspace{-.25cm}
    \end{enumerate}
    Let $G_a$ be the graph with vertex set $V_a$ (whose elements we call \emph{part-vertices}) and an edge between two part-vertices $X\subseteq V(G)$ and $Y\subseteq V(G)$ if there is an edge in $G$ between some $x\in X$ and $y\in Y$. Note that there are at most $|B_a|+3$ part-vertices in $G_a$; on for each each $B_a$ and one vertex for each neighbour of $a$.

    Now let $t\in V(T)$ and $e\in E(T)$ be an incident edge. Observe that for each part-vertex $X\in V_t$, there is a unique part-vertex $Y\in V_e$ so that $X\subseteq Y$. Let $g_t^e\colon V_t\to V_e$ denote the map defined this way, and observe that it defines an epimorphism (surjective homomorphism) $G_t\to G_e$.

    We will show that the limit of $\Ta$ is $G$. To that end, for each $a\in V(T)\cup E(T)$, let $g_a\colon G\to G_a$ be the epimorphism sending each vertex to the part-vertex containing it, and observe that $(G,\{g_a\})$ is a cone over $\Ta$ with $G$ as its apex. It remains to prove that $G$ is universal.
    
    Let $(H,\{h_a\})$ be another cone over $\Ta$. Let $u\in V(H)$. We make the following claims:

    \textbf{Claim.} There exists a node $t\in V(T)$ for which $h_t(u)$ is a singleton.
    
    \textit{Proof of claim.} Let $t\in V(T)$ be arbitrary. If $h_t(u)$ is a singleton, then we are done, so suppose this is not the case. Then $h_t(u)=P_C$ for some connected component $C$ of $T-t$. Let $tt'$ be the unique edge of $T$ incident with $C$ but not contained within it. Then $P_C$ is also one of the parts of $G_{tt'}$ and $h_{tt'}(u)=P_C$ by the commutativity of the cone $(H,\{h_a\})$. Since otherwise we would be done, assume that $h_{t'}(u)$ is not a singleton. Then $h_{t'}(u)=P_D$ for some $D\subseteq C$ a connected component of $T-t'$. Starting again with $P_D$ in place of $P_C$, we iterate. Since $T$ is finite, this process must terminate in finding our required singleton.
    \qed

    \textbf{Claim.} There exists a unique vertex $x\in V(G)$ so that $x\in h_a(u)$ for all $a\in V(T)\cup E(T)$.
    
    \textit{Proof of claim.} By the first claim, $h_a(u)$ is a singleton $\{x\}$ for some $a\in V(T)$ and $x\in V(G)$. Observe that by the commutativity of the cones, $h_a(u)=\{x\}$ for each $a$ in the contiguous subtree $T[x]$ of $T$ whose bags contain $x$. For any $a\in V(T)\cup E(T)$ not in $T[x]$, it is clear by induction that $h_a(u)$ is the unique part representing the connected component of $T-a$ containing $T[x]$. That is, the part of $V_a$ containing $x$. Hence the choice of $x$ satisfies our claim.
    \qed

    Observe that the map that sends each $u\in V(H)$ to $x\in V(G)$ gives a function $f\colon H\to G$ that commutes with both cones, and that furthermore this function is unique with this property.
    It remains to show that $f$ is a graph homomorphism. So let $(u,v)$ be an edge in $H$. Observe that there is an edge between from $h_a(u)=g_a(f(u))$ to $h_a(v)g_a(f(v))$ in $G_a$ for each $a\in V(T)\cup E(T)$. Hence $(f(u),f(v))$ is an edge in $G$. Hence $f$ is a homomorphism and we conclude that $G$ is the limit of $\Ta$.

    Our proof is completed by observing that for every $a\in V(T)\cup E(T)$, the forest $T-a$ has at most three connected components, and hence the graph $G_a$ has at most $|B_a|+3\le w-1+3=w+2$ vertices.
\end{proof}